\begin{HVTPage}{Capítulo 3}{Pirólisis de residuos}{Ruta termoquímica}
\HVTLead{La pirólisis transforma residuos orgánicos mediante calentamiento en ausencia o fuerte limitación de oxígeno. Sus productos principales son biocarbón, bioaceite y gas; la obtención de una corriente rica en H2 requiere reformado y acondicionamiento posteriores.}
\HVTEquation{\text{Residuos orgánicos}+\text{calor}\rightarrow\text{biocarbón}+\text{bioaceite}+\text{gas}}{}
\HVTText{La expresión es una representación conceptual de productos y no una ecuación estequiométrica. Sus fracciones dependen de la biomasa, la temperatura, la velocidad de calentamiento y el tiempo de residencia. \HVTCite{3}}
\begin{HVTFlow}\HVTFlowItem{1}{Preparación}{Clasificación, secado, reducción de tamaño y caracterización.}\HVTFlowItem{2}{Pirólisis}{Control de temperatura, tiempo de residencia y atmósfera.}\HVTFlowItem{3}{Conversión del gas}{Reformado, conversión de CO y purificación de H2.}\end{HVTFlow}
\end{HVTPage}

\begin{HVTPage}{Capítulo 3}{Pirólisis de residuos}{Variables y precisión técnica}
\begin{HVTTable}{Variable}{Símbolo}{Unidad}
\HVTTableRow{Humedad del residuo}{$X_{H_2O}$}{\% masa}
\HVTTableRow{Temperatura del reactor}{$T_p$}{°C}
\HVTTableRow{Tiempo de residencia}{$t_r$}{s o min}
\HVTTableRow{Rendimiento de gas}{$Y_g$}{kg gas/kg residuo}
\HVTTableRow{Fracción de H2 en gas}{$x_{H_2,g}$}{\% vol.}
\end{HVTTable}
\begin{HVTWarning}{Precisión técnica}El desempeño depende de la composición del residuo, especialmente humedad, cenizas y fracciones de celulosa, hemicelulosa y lignina. \HVTCite{4} La producción de H2 a partir de bioaceite requiere etapas posteriores, como reformado catalítico con vapor y purificación. \HVTCite{5} No debe asignarse un rendimiento único sin caracterización y ensayos.\end{HVTWarning}
\end{HVTPage}
